\documentclass[12pt,a4paper]{report}
\usepackage{fontspec}
\IfFontExistsTF{TeX Gyre Termes}{\setmainfont{TeX Gyre Termes}}{\setmainfont{Liberation Serif}}
\IfFontExistsTF{TeX Gyre Heros}{\setsansfont{TeX Gyre Heros}}{\setsansfont{Liberation Sans}}
\setmonofont[Scale=0.84]{DejaVu Sans Mono}
\usepackage[vietnamese]{babel}
\usepackage[a4paper,margin=1in]{geometry}
\usepackage{amsmath,amssymb,mathtools,bm}
\usepackage{booktabs,tabularx,array,longtable}
\usepackage{graphicx,float}
\usepackage[hidelinks]{hyperref}
\usepackage{fancyhdr,setspace,titlesec,caption,microtype,needspace}
\usepackage{enumitem}
\setstretch{1.15}
\setlength{\parskip}{0.3em}
\setlength{\parindent}{0pt}
\setlength{\headheight}{16pt}
\setcounter{tocdepth}{1}
\setcounter{secnumdepth}{2}
\captionsetup{font=small,labelfont=bf}
\numberwithin{equation}{chapter}
\titleformat{\chapter}[display]{\normalfont\bfseries}{\filleft\Large Chương \thechapter}{0.3em}{\filright\huge}[\vspace{0.4em}\titlerule]
\titlespacing*{\chapter}{0pt}{-10pt}{20pt}
\titleformat{\section}{\large\bfseries}{\thesection}{0.6em}{}
\titleformat{\subsection}{\normalsize\bfseries}{\thesubsection}{0.6em}{}
\pagestyle{fancy}
\fancyhf{}
\lhead{\small CO5085 · Bài tập A1}
\rhead{\small Huỳnh Thành Đạt · 2570161}
\cfoot{\thepage}
\fancypagestyle{plain}{\fancyhf{}\lhead{\small CO5085 · Bài tập A1}\rhead{\small Huỳnh Thành Đạt · 2570161}\cfoot{\thepage}}
\newcommand{\R}{\mathbb{R}}
\newcommand{\norm}[1]{\left\lVert #1 \right\rVert}
\newcommand{\code}[1]{\texttt{\detokenize{#1}}}
\newcommand{\resultfigure}[3]{\IfFileExists{generated/figures/#1.png}{\begin{figure}[htbp]\centering\includegraphics[width=#3\textwidth]{generated/figures/#1.png}\caption{#2}\end{figure}}{}}
\hypersetup{pdftitle={CO5085 - Bao cao A1},pdfauthor={Huynh Thanh Dat - 2570161}}
\emergencystretch=2em
\begin{document}
\hypersetup{pageanchor=false}
\begin{titlepage}
\centering
{\large ĐẠI HỌC QUỐC GIA THÀNH PHỐ HỒ CHÍ MINH\par}
{\large\bfseries TRƯỜNG ĐẠI HỌC BÁCH KHOA\par}
{\large KHOA KHOA HỌC VÀ KỸ THUẬT MÁY TÍNH\par}
\vspace{1.2cm}
\includegraphics[width=3cm]{assets/hcmut-logo.png}\par
\vspace{0.8cm}
{\large\bfseries HỌC SÂU VÀ ỨNG DỤNG\par}
{\large\bfseries TRONG THỊ GIÁC MÁY TÍNH\par}
{\large Mã môn học: CO5085\par}
\vspace{0.9cm}
\rule{\textwidth}{0.6pt}\par\vspace{0.35cm}
{\Huge\bfseries BÁO CÁO BÀI TẬP A1\par}
\vspace{0.45cm}
{\Large CNN vs Transformer trên CIFAR-100\par}
{\Large ResNet18 và Swin-T\par}
\vspace{0.35cm}\rule{\textwidth}{0.6pt}\par
\vspace{0.7cm}
{\large Học kỳ 261 · Năm học 2026–2027\par}
\vspace{0.7cm}
\begin{tabular}{@{}ll@{}}
Giảng viên: & Lê Thành Sách\\
& Nguyễn Quốc Minh\\[0.5em]
Sinh viên: & Huỳnh Thành Đạt\\
Mã số sinh viên: & 2570161\\
Hình thức: & Cá nhân
\end{tabular}
\vfill
{\small Thành phố Hồ Chí Minh, tháng 10 năm 2026\par}
\end{titlepage}
\hypersetup{pageanchor=true}
\pagenumbering{roman}
{\setstretch{1.0}\small\tableofcontents}
\clearpage
\pagenumbering{arabic}

\chapter{Bài toán và dữ liệu}
\section{Mục tiêu và phạm vi}
A1 so sánh CNN và Transformer trên CIFAR-100, tập trung vào ảnh hưởng của
khởi tạo pretrained và mức độ đóng băng backbone. Hai kiến trúc được chọn là
ResNet18 và Swin-T của torchvision. Mỗi kiến trúc có bốn chế độ huấn luyện,
chạy với ba seed 42, 43 và 44, tổng cộng 24 lượt chạy.

Các so sánh chính gồm huấn luyện từ đầu với pretrained, và head-only với
partial/full fine-tuning. Accuracy, macro-F1, thời gian và bộ nhớ được dùng
để đánh giá hiệu quả cùng chi phí của từng chế độ.

\section{CIFAR-100 và tình trạng license}
CIFAR-100 gồm 60.000 ảnh màu $32\times32$, 100 nhãn chi tiết (fine classes)
và 20 nhóm coarse classes. Bài này dự đoán toàn bộ 100 fine classes; không
chọn subset ảnh hoặc lớp. Mỗi lớp có 500 ảnh official train và 100 ảnh
official test \cite{cifar}. Số lớp lớn hơn bài E1 và ảnh độ phân giải thấp
có thể làm những lớp gần nhau khó phân biệt.

Nguồn chính là trang CIFAR của Toronto, của Alex Krizhevsky, Vinod Nair và
Geoffrey Hinton. Trang này không nêu license tường minh; thẻ dataset chính
thức \code{uoft-cs/cifar100} trên Hugging Face ghi \emph{license: unknown}
\cite{cifar,hf}. Vì vậy, tình trạng license của dữ liệu được ghi là chưa rõ.
Ảnh gốc được tải từ Toronto để thực hiện thí nghiệm và lưu ngoài repository.

\section{Split và tiền xử lý}
Từ 50.000 ảnh train gốc, shuffle indices riêng từng lớp bằng NumPy
\code{default_rng(2026)}. Lấy 50 ảnh/lớp làm validation, 450 ảnh/lớp làm train:
45.000 train và 5.000 validation. Indices được sắp tăng sau khi chọn và lưu
cùng hash. Official test 10.000 ảnh nằm riêng, không dùng cập nhật tham số
hoặc chọn checkpoint. Train seed 42/43/44 không thay split seed 2026.

Train dùng RandomCrop $32\times32$ với padding 4, horizontal flip xác suất 0,5,
sau đó bilinear resize $96\times96$ có antialias. Val/test chỉ resize, không
crop ngẫu nhiên hoặc flip. Cả hai họ và mọi chế độ dùng ImageNet mean
$(0{,}485,0{,}456,0{,}406)$, std $(0{,}229,0{,}224,0{,}225)$ sau ToTensor.
Không dùng trực tiếp recipe \code{weights.transforms()} của pretrained.

Augmentation dùng seed theo train seed, epoch và chỉ số ảnh để giữ phép
biến đổi nhất quán khi resume.

\chapter{Kiến trúc và chế độ huấn luyện}
\section{ResNet18: đối chứng CNN}
Giữ kiến trúc ResNet18 chuẩn torchvision, gồm stem convolution $7\times7$,
max pooling và bốn residual stages. Một residual block cập nhật
\begin{equation}
h_{l+1}=\operatorname{ReLU}\left(F_l(h_l)+S_l(h_l)\right),
\end{equation}
trong đó $S_l$ là identity hoặc projection khi đổi số kênh/stride. Kernel
convolution chia sẻ trọng số theo vị trí và khai thác quan hệ cục bộ. Sau
adaptive average pooling, head Linear được thay để dự đoán 100 lớp; không
thay stem thành biến thể ResNet chuyên cho CIFAR.

Bản pretrained dùng \code{ResNet18_Weights.IMAGENET1K_V1}. Ba chế độ là
chỉ học \code{fc}; học \code{layer4} và \code{fc}; học toàn bộ. Bản scratch
khởi tạo mới và học toàn bộ. BatchNorm trong phần frozen ở eval, nên cả
tham số lẫn running mean/variance và bộ đếm batch được giữ nguyên.

\section{Swin-T: đối chứng Transformer}
Swin-T dùng patch embedding stride4, các stage Transformer theo cấp độ và
patch merging. Attention trên cửa sổ có dạng
\begin{equation}
\operatorname{Attention}(Q,K,V)
=\operatorname{softmax}\left(\frac{QK^\top}{\sqrt{d}}+B+M\right)V,
\end{equation}
trong đó $B$ là relative position bias, $M$ là mask của shifted windows.
Dịch cửa sổ giúp trao đổi thông tin giữa cửa sổ qua các block; việc giới hạn
attention trong cửa sổ khác với global attention của ViT cổ điển.

Swin-T giữ kiến trúc chuẩn torchvision, với \code{head} được thay bằng
Linear 100 lớp. Đầu vào $96\times96$ tạo lưới patch đầu $24\times24$;
các stage muộn có lưới nhỏ hơn. Kích thước cửa sổ mặc định là 7, với padding
nội bộ ở những stage có kích thước không chia hết cho cửa sổ.

Bản pretrained dùng \code{Swin_T_Weights.IMAGENET1K_V1}. Chế độ partial chỉ
học stage cuối \code{features.7}, \code{norm} cuối và \code{head}; patch
merging đứng trước stage cuối vẫn frozen. Head-only chỉ học head. Frozen
stages được giữ ở eval, bao gồm stochastic depth, để đặc trưng của phần
đóng băng ổn định trong quá trình học head. Full và scratch cập nhật toàn bộ mô hình.

\section{Ma trận đối chứng}
\begin{table}[htbp]\centering
\begin{tabular}{lll}
\toprule
Chế độ & Khởi tạo backbone & Phần được cập nhật\\\midrule
Scratch & Ngẫu nhiên & Toàn bộ\\
Head-only & ImageNet pretrained & Head100\\
Partial & ImageNet pretrained & Stage cuối và head; Swin thêm norm\\
Full & ImageNet pretrained & Toàn bộ\\\bottomrule
\end{tabular}
\caption{Bốn chế độ cho mỗi họ kiến trúc, cùng split và ba seed.}
\end{table}
Scratch cập nhật toàn bộ mô hình từ khởi tạo ngẫu nhiên. Số tham số tổng và
trainable được đếm trực tiếp từ từng mô hình và lưu trong metadata.
ResNet18 và Swin-T khác nhau về kiến trúc, capacity và recipe pretrained;
các kết quả phản ánh tổng thể những khác biệt này trong cùng thiết lập.

\chapter{Thiết lập thực nghiệm và phép đo}
\section{Thiết lập huấn luyện}
Mỗi lượt huấn luyện có 20 epoch, batch 128 (giữ batch cuối), vòng lặp huấn luyện trực tiếp,
AdamW với weight decay $10^{-4}$, $\beta=(0{,}9,0{,}999)$, epsilon $10^{-8}$.
Learning rate ban đầu là $10^{-3}$ cho scratch/head và $10^{-4}$ cho
pretrained partial/full. Cosine scheduler bước một lần sau epoch:
\begin{equation}
\eta_e=\frac{\eta_0}{2}\left[1+\cos\left(\frac{\pi e}{20}\right)\right],
\end{equation}
với $e=0$ trước epoch đầu. Loss cross-entropy không label smoothing:
\begin{equation}
\mathcal{L}=-\frac1B\sum_{i=1}^{B}\log
\frac{\exp z_{i,y_i}}{\sum_{c=1}^{100}\exp z_{i,c}}.
\end{equation}
Clip global gradient norm ở 1 trước optimizer.step. BF16 autocast áp dụng
cho tất cả chế độ; tham số và Adam moments giữ FP32, không GradScaler,
TF32 tắt và không dùng compile. Các phép đo được thiết lập trên RTX 5090 với
8 workers và 4 CPU threads.

Ngân sách 20 epoch được giữ cố định giữa các chế độ. Learning rate của
head/scratch và partial/full được đặt khác nhau để phù hợp với mức độ cập
nhật backbone. Kết quả vì vậy được phân tích theo từng thiết lập đầy đủ;
mức hội tụ còn được đối chiếu qua đường cong train/validation.

\section{Chọn checkpoint và mở test}
Best chọn bằng val accuracy lớn nhất; hòa accuracy thì lấy CE nhỏ nhất,
hòa cả hai giữ epoch trước. Chỉ val được dùng ra quyết định trong train.
Test được đánh giá ở một bước riêng sau khi hoàn tất các lượt huấn luyện.
Các cấu hình, learning rate và seed được cố định trước bước này.

Accuracy và macroF1 đều tính trên 100 lớp. Với mỗi lớp $c$,
\begin{equation}
F1_c=\frac{2TP_c}{2TP_c+FP_c+FN_c},\qquad
F1_{macro}=\frac1{100}\sum_{c=1}^{100}F1_c.
\end{equation}
Trường hợp mẫu số bằng0 có F1 bằng0. Confusion matrix, support, precision,
recall, F1 theo lớp và từng prediction được lưu để kiểm tra. Gallery lấy tối
đa 8 lỗi confidence cao nhất của seed 42.

\section{Chi phí và môi trường}
Training time bao gồm data loading/augmentation và train pass, có CUDA
synchronize trước/sau. Epoch time thêm validation. Model construction,
checkpoint và report export nằm ngoài hai khoảng đo; không gọi training time
là tổng wall time của job. Peak GPU memory là peak allocated trong train và
validation của mỗi epoch, không phải reserved memory hoặc VRAM toàn hệ thống.

Metadata lưu GPU, torch/torchvision, CUDA/cuDNN, workers, CPU threads và
chế độ deterministic. Các phép đo thời gian của cùng cấu hình được gộp khi
những thông số môi trường này thống nhất.

\section{Checkpoint và thông tin tái lập}
Run lưu \code{last.pt} với model, optimizer, scheduler, RNG Python/NumPy/
PyTorch/CUDA, loader generators và history; \code{best.pt} chỉ lưu model và
selection. Không giữ best optimizer hay checkpoint mỗi epoch. Atomic writes
và khóa run bảo vệ các tệp kết quả. Resume từ epoch hoàn tất gần nhất, không giữa
batch; best write bị ngắt có thể khôi phục từ last khi last là best epoch.

Config/source/data/split/environment hashes bảo vệ resume và test. Frozen
parameters/buffers được fingerprint kiểm tra sau epoch. Git commit/dirty
được ghi riêng để đối chiếu phiên bản mã nguồn. Checkpoint được lưu ngoài
Git; JSON/CSV và checksums được giữ trong generated/evidence. Sai khác
GPU hoặc kernel có thể ảnh hưởng khả năng tái lập từng bit.

\section{Mã nguồn và kiểm thử}
Mã A1 nằm trong thư mục \code{src/A1/}, cấu hình ở \code{configs/a1_experiments.json}.
Thí nghiệm được chạy bằng \code{A1.run_suite} trên CUDA với 8 workers;
\code{--resume} khôi phục lượt chạy bị gián đoạn. \code{A1.export_report}
tổng hợp bảng và hình; \code{scripts.build_report} biên dịch LaTeX thành PDF.
Các lệnh đầy đủ được ghi trong \code{src/A1/README.md}.

Các kiểm thử số học và protocol không tải dữ liệu hoặc trọng số pretrained. Phạm vi gồm split theo chỉ số, accuracy/F1, đóng băng tham số và buffers, augmentation theo epoch/worker, khôi phục checkpoint/RNG và kiểm tra signature.

\chapter{Kết quả thực nghiệm và thảo luận}
\section{So sánh các chế độ huấn luyện}
\begin{table}[htbp]\centering\small
\setlength{\tabcolsep}{5pt}
\begin{tabular}{llrrr}
\toprule
Kiến trúc & Chế độ & Val (\%) & Test (\%) & Macro-F1 \\
\midrule
ResNet18 & Scratch & $63.83 \pm 0.51$ & $64.35 \pm 0.11$ & $0.6422 \pm 0.0012$ \\
ResNet18 & Head-only & $61.89 \pm 0.24$ & $60.95 \pm 0.29$ & $0.6081 \pm 0.0030$ \\
ResNet18 & Partial & $74.50 \pm 0.31$ & $73.96 \pm 0.12$ & $0.7396 \pm 0.0011$ \\
ResNet18 & Full & $79.29 \pm 0.10$ & $78.97 \pm 0.04$ & $0.7898 \pm 0.0006$ \\
Swin-T & Scratch & $2.81 \pm 1.40$ & $2.76 \pm 1.26$ & $0.0093 \pm 0.0069$ \\
Swin-T & Head-only & $72.45 \pm 0.10$ & $71.67 \pm 0.09$ & $0.7154 \pm 0.0009$ \\
Swin-T & Partial & $78.99 \pm 0.31$ & $78.25 \pm 0.22$ & $0.7818 \pm 0.0021$ \\
Swin-T & Full & $86.89 \pm 0.02$ & $86.63 \pm 0.22$ & $0.8663 \pm 0.0020$ \\
\bottomrule
\end{tabular}

\caption{Accuracy và macro-F1 của tám cấu hình, trung bình và SD mẫu qua ba seed. Accuracy theo phần trăm, F1 theo thang 0--1.}
\end{table}
Kết quả được tổng hợp bằng trung bình và độ lệch chuẩn mẫu qua các seed.
Danh sách lượt chạy và checkpoint tương ứng được lưu trong manifest.

\section{Thời gian, bộ nhớ và số tham số}
\begin{table}[htbp]\centering\small
\setlength{\tabcolsep}{4pt}
\begin{tabular}{llrrrr}
\toprule
Kiến trúc & Chế độ & Tổng (M) & Học (M) & Train (s) & VRAM (MiB) \\
\midrule
ResNet18 & Scratch & 11.228 & 11.228 & $253.4 \pm 1.2$ & 522.5 \\
ResNet18 & Head-only & 11.228 & 0.051 & $155.9 \pm 2.9$ & 252.1 \\
ResNet18 & Partial & 11.228 & 8.445 & $194.1 \pm 0.2$ & 347.3 \\
ResNet18 & Full & 11.228 & 11.228 & $257.6 \pm 3.2$ & 522.4 \\
Swin-T & Scratch & 27.596 & 27.596 & $727.1 \pm 4.8$ & 2456.9 \\
Swin-T & Head-only & 27.596 & 0.077 & $265.3 \pm 2.1$ & 502.2 \\
Swin-T & Partial & 27.596 & 14.262 & $347.5 \pm 0.4$ & 663.6 \\
Swin-T & Full & 27.596 & 27.596 & $724.4 \pm 1.7$ & 2456.9 \\
\bottomrule
\end{tabular}

\caption{Tổng tham số và tham số được học theo triệu (M); thời gian train đủ 20 epoch; peak allocated memory trung bình theo MiB. Phạm vi phép đo ở Mục 3.3.}
\end{table}

\section{Thảo luận kết quả}
Kết quả gồm 24 lượt chạy, với ba seed 42, 43 và 44 cho mỗi cấu hình. Chênh lệch accuracy được tính theo cặp cùng seed; các khoảng $\pm$ là độ lệch chuẩn mẫu, không phải khoảng tin cậy.

\subsection{Khởi tạo và khả năng thích nghi biểu diễn}
ResNet18 head-only so với scratch có chênh lệch -3.40 $\pm$ 0.33 điểm phần trăm. Backbone đóng băng chỉ cho phép head tuyến tính điều chỉnh ranh giới phân loại trên đặc trưng có sẵn. Kết quả cho thấy đặc trưng ImageNet cố định chưa đủ phù hợp với thiết lập CIFAR-100 này. Mở stage cuối tăng +9.61 $\pm$ 0.09 điểm phần trăm so với scratch; mở toàn bộ tăng +14.62 $\pm$ 0.11 điểm phần trăm. Khả năng cập nhật biểu diễn đóng vai trò đáng kể đối với ResNet18.

Swin-T scratch chỉ đạt $2.76 \pm 1.26$\% trên test. Ở epoch 20, seed 43 và 44 có validation accuracy 1\%, train accuracy xấp xỉ 1\% và CE gần $\ln(100)\simeq4{,}605$. Cả train lẫn validation đều thấp, cho thấy mô hình học rất ít trong thiết lập này. Seed 42 cải thiện hơn nhưng vẫn xa các bản pretrained. Ngân sách 20 epoch, learning rate $10^{-3}$ và lịch cosine không warm-up chưa tạo được một đối chứng scratch hội tụ tốt. Thực nghiệm không tách riêng nguyên nhân của từng yếu tố và không khảo sát thêm siêu tham số sau khi mở test.

Swin-T head-only đạt $71.67 \pm 0.09$\%, partial đạt $78.25 \pm 0.22$\%, và full đạt $86.63 \pm 0.22$\%. Full cao hơn scratch +83.87 $\pm$ 1.40 điểm phần trăm; mức chênh lệch này phản ánh cả lợi ích của pretrained lẫn đối chứng scratch chưa học tốt. Kết quả không đủ để kết luận Transformer luôn phụ thuộc pretrained hoặc luôn vượt CNN trong những ngân sách khác.

\subsection{Accuracy và chi phí huấn luyện}
Khi cùng fine-tuning toàn bộ, Swin-T cao hơn ResNet18 +7.66 $\pm$ 0.18 điểm phần trăm, với macro-F1 trung bình 0.8663 so với 0.7898. Thời gian train tăng 2.81 lần và peak allocated memory tăng 4.70 lần. Đây là sự đánh đổi giữa accuracy và chi phí trong cùng ngân sách epoch; hai kiến trúc vẫn khác số tham số và quá trình pretraining.

Head-only giảm peak allocated memory so với full từ 522.4 xuống 252.1 MiB ở ResNet18, và từ 2456.9 xuống 502.2 MiB ở Swin-T. Đóng băng backbone giảm bộ nhớ gradient và trạng thái optimizer, nhưng vẫn cần bước forward qua toàn bộ backbone.

\Needspace{7\baselineskip}
\subsection{Đường cong học và lỗi phân loại}
Ở epoch 20, train/validation accuracy của các bản full có khoảng cách rõ: ResNet18 99.07\%/79.18\%; Swin-T 98.68\%/86.75\%. Train cải thiện nhanh hơn validation, trong khi validation CE còn cao hơn train CE; đây là dấu hiệu overfit. Điểm test trong bảng lấy từ checkpoint tốt nhất theo validation, không nhất thiết là epoch 20.

ResNet18 full ở seed 42: lỗi nhiều nhất theo một chiều là oak\_tree thành maple\_tree, 23 ảnh. Lớp oak\_tree có 100 ảnh trong tập test.

Swin-T full ở seed 42: lỗi nhiều nhất theo một chiều là oak\_tree thành maple\_tree, 22 ảnh. Lớp oak\_tree có 100 ảnh trong tập test.

Hai mô hình cùng nhầm oak\_tree với maple\_tree, gợi ý khó khăn khi phân biệt các lớp cây có hình dáng và kết cấu gần nhau trong ảnh $32\times32$. Đây là cách diễn giải từ ma trận lỗi và các ví dụ quan sát, không phải phép đo riêng về mức tương đồng thị giác. Những lỗi confidence cao cho thấy accuracy cao vẫn đi kèm một số dự đoán sai rất tự tin; calibration chưa được đánh giá riêng.


\section{Đường cong và lỗi phân loại}
Các hình dùng seed 42 đã định trước, không chọn seed theo kết quả test.
Ma trận lỗi dùng chỉ số lớp 0--99; danh sách tên lớp và số đếm đầy đủ nằm
trong các tệp CSV kèm báo cáo.
\resultfigure{resnet18_scratch_curves}{ResNet18 scratch: train/val CE và accuracy ở seed 42.}{0.95}
\resultfigure{swin_t_scratch_curves}{Swin-T scratch: train/val CE và accuracy ở seed 42.}{0.95}
\resultfigure{resnet18_full_curves}{ResNet18 pretrained full: đường cong ở seed 42.}{0.95}
\resultfigure{swin_t_full_curves}{Swin-T pretrained full: đường cong ở seed 42.}{0.95}
\resultfigure{resnet18_full_confusion_test}{Confusion matrix test của ResNet18 full ở seed 42; hàng là nhãn thực tế.}{0.9}
\resultfigure{swin_t_full_confusion_test}{Confusion matrix test của Swin-T full ở seed 42.}{0.9}
\resultfigure{resnet18_full_errors_test}{Các dự đoán sai có confidence cao của ResNet18 full ở seed 42, hiển thị ảnh gốc $32\times32$.}{0.95}
\resultfigure{swin_t_full_errors_test}{Các dự đoán sai có confidence cao của Swin-T full ở seed 42.}{0.95}

\chapter{Kết luận và giới hạn}
\section{Kết luận}
Trong thiết lập CIFAR-100 đầy đủ với 20 epoch, Swin-T fine-tuning toàn bộ cho test accuracy cao nhất, đạt $86.63 \pm 0.22$\% và macro-F1 0.8663. Fine-tuning stage cuối hoặc toàn bộ cải thiện kết quả so với chỉ học head ở cả hai kiến trúc. ResNet18 scratch học được biểu diễn hữu ích trong ngân sách này, còn Swin-T scratch chưa học tốt. Lợi thế accuracy của Swin-T full đi kèm thời gian và bộ nhớ cao hơn ResNet18 full.


\section{Giới hạn}
Ảnh gốc $32\times32$ và đầu vào $96\times96$ khác với kích thước phổ biến
$224\times224$ của pretrained. Thực nghiệm chỉ khảo sát CIFAR-100, hai kiến
trúc và ngân sách 20 epoch; đặc biệt, scratch có thể tiếp tục cải thiện nếu
được huấn luyện lâu hơn. Hai họ còn khác capacity và recipe ImageNet.
Ba seed mô tả biến thiên giữa các lượt khởi tạo, với cùng phần cứng và
thiết lập workers. Tình trạng license của dữ liệu vẫn là chưa rõ.

\subsection{Ảnh trùng giữa các split}
Audit toàn bộ pixel và nhãn, khớp hash dữ liệu gốc, phát hiện 2 nội dung ảnh trùng train--validation và 10 trùng train--test; không có trùng validation--test. Split và checkpoint gốc được giữ nguyên. Loại 10 ảnh test trùng train khỏi prediction đã lưu cho accuracy trung bình 79.009\% với ResNet18 full và 86.677\% với Swin-T full. Đây là kiểm tra độ nhạy, không phải train lại sau loại trùng. Hai nội dung trùng train--validation vẫn có thể ảnh hưởng chọn checkpoint; near duplicates chưa được đánh giá. Chi tiết nằm trong \code{docs/training-evidence/cifar100-pixel-audit.json} và các tệp overlap/sensitivity.

Báo cáo sử dụng AI được đính kèm tại tệp \code{AI_USE_REPORT.md} trong repository.

\begin{thebibliography}{9}
\addcontentsline{toc}{chapter}{Tài liệu tham khảo}
\bibitem{assignment} Lê Thành Sách, Nguyễn Quốc Minh. \emph{Bài tập lớn A1: CNN vs Transformer trên tập ảnh lớn}. CO5085, học kỳ 261.
\bibitem{cifar} A. Krizhevsky. \emph{Learning Multiple Layers of Features from Tiny Images}. 2009. Trang dataset: \url{https://www.cs.toronto.edu/~kriz/cifar.html}.
\bibitem{hf} uoft-cs. \emph{CIFAR100 dataset card}; license: unknown. \url{https://huggingface.co/datasets/uoft-cs/cifar100}.
\bibitem{resnet} K. He và cộng sự. \emph{Deep Residual Learning for Image Recognition}. CVPR2016. \url{https://arxiv.org/abs/1512.03385}.
\bibitem{swin} Z. Liu và cộng sự. \emph{Swin Transformer: Hierarchical Vision Transformer using Shifted Windows}. ICCV2021. \url{https://arxiv.org/abs/2103.14030}.
\bibitem{torchvision} Torchvision. ResNet18 và Swin-T model API; sử dụng weight enum IMAGENET1K\_V1. \url{https://docs.pytorch.org/vision/stable/models.html}.
\end{thebibliography}
\end{document}
